\documentclass[12pt]{article}
\usepackage{hyperref}
\usepackage{listings}
\usepackage[margin=1in]{geometry}
\usepackage{enumitem}
\usepackage{multicol}
\usepackage{array}
\usepackage{titlesec}
\usepackage{helvet}
\renewcommand{\familydefault}{\sfdefault}
\usepackage{amsmath}     % For math equations
\usepackage{amssymb}     % For advanced math symbols
\usepackage{amsfonts} % For math fonts
\usepackage{gvv}
\usepackage{esint}
\usepackage[utf8]{inputenc}
\usepackage{graphicx}
\usepackage{pgfplots}
\pgfplotsset{compat=1.18}
\titleformat{\section}{\bfseries\large}{\thesection.}{1em}{}
\setlength{\parindent}{0pt}
\setlength{\parskip}{6pt}
\usepackage{multirow}
\usepackage{float}
\usepackage{caption}

\begin{document}

\textbf{Problem 12.526}

Let $A$ be a $7\times7$ matrix such that 
\begin{align}
2A^2 - A^4 = I_7
\end{align}
where $I_7$ is the identity matrix. If $A$ has two distinct eigenvalues and each eigenvalue has geometric multiplicity $3$, then the total number of nonzero entries in the Jordan canonical form of $A$ equals:
\begin{align*}
\text{(A) } 6 \qquad \text{(B) } 7 \qquad \text{(C) } 8 \qquad \text{(D) } 9
\end{align*}

\vspace{1em}

\section*{Given Data}

\begin{table}[H]
\centering
\renewcommand{\arraystretch}{1.4}
\begin{tabular}{|c|c|c|}
\hline
\textbf{Symbol} & \textbf{Meaning} & \textbf{Value} \\
\hline
$A$ & Square matrix & $7\times7$ \\
\hline
Equation & Polynomial relation & $2A^2 - A^4 = I_7$ \\
\hline
Eigenvalues & Distinct eigenvalues of $A$ & Two distinct ($\lambda = \pm1$) \\
\hline
Geometric multiplicity (GM) & $\dim(\ker(A-\lambda I))$ & $3$ for each eigenvalue \\
\hline
Total GM & Sum over both eigenvalues & $3 + 3 = 6$ \\
\hline
\end{tabular}
\end{table}

\vspace{1em}

\section*{Theory}

\subsection*{1. Geometric Multiplicity}

For a square matrix $A$ and eigenvalue $\lambda$,
\begin{align}
\text{GM}(\lambda) = \dim(\ker(A - \lambda I)).
\end{align}
It represents the number of linearly independent eigenvectors for $\lambda$.  

If the \textbf{algebraic multiplicity (AM)} (how many times $\lambda$ repeats in the characteristic polynomial) equals the GM, then $A$ is diagonalizable.  
Otherwise, $A$ is \textbf{defective} and has Jordan blocks larger than $1\times1$.

\subsection*{2. Jordan Canonical Form}

Every square matrix $A$ is similar to a matrix $J$ such that
\begin{align}
A = PJP^{-1}
\end{align}
where $J$ is the \textbf{Jordan canonical form} and $P$ is an invertible matrix whose columns are eigenvectors and generalized eigenvectors.

A Jordan block of size $k$ corresponding to an eigenvalue $\lambda$ is:
\begin{align}
J_k(\lambda) =
\myvec{
\lambda & 1 & 0 & \cdots & 0\\
0 & \lambda & 1 & \cdots & 0\\
\vdots & & \ddots & \ddots & \vdots\\
0 & \cdots & 0 & \lambda & 1\\
0 & \cdots & \cdots & 0 & \lambda
}.
\end{align}

The $1$ above the diagonal connects a \textbf{generalized eigenvector} to a true eigenvector.  
For example:
\begin{align}
\myvec{\lambda & 1\\0 & \lambda}
\end{align}
represents one eigenvector $\vec{v_1}$ satisfying $(A-\lambda I)\vec{v_1}=0$, and one generalized eigenvector $\vec{v_2}$ satisfying $(A-\lambda I)\vec{v_2}=\vec{v_1}$.

\subsection*{3. How Jordan Block Sizes are Determined Using Null Spaces}

We define
\begin{align}
N_k(\lambda) = \ker((A - \lambda I)^k).
\end{align}

The spaces satisfy
\begin{align}
N_1(\lambda) \subseteq N_2(\lambda) \subseteq N_3(\lambda) \subseteq \cdots
\end{align}
and their dimensions reveal Jordan block structure:

\begin{table}[H]
\centering
\renewcommand{\arraystretch}{1.3}
\begin{tabular}{|c|c|}
\hline
\textbf{Quantity} & \textbf{Meaning} \\
\hline
$\dim N_1(\lambda)$ & Number of Jordan blocks (= GM) \\
\hline
$\dim N_2(\lambda) - \dim N_1(\lambda)$ & Number of blocks of size $\ge 2$ \\
\hline
$\dim N_3(\lambda) - \dim N_2(\lambda)$ & Number of blocks of size $\ge 3$ \\
\hline
\end{tabular}
\end{table}

Thus, computing $\ker((A-\lambda I)^k)$ for successive $k$ directly determines how many blocks of each size exist.

\section*{Solution}

\subsection*{Step 1: Simplify the given polynomial}
\begin{align}
2A^2 - A^4 = I_7
&\Rightarrow A^4 - 2A^2 + I = 0\\
&\Rightarrow (A^2 - I)^2 = 0
\end{align}

\subsection*{Step 2: Find eigenvalues}
Let $\lambda$ be an eigenvalue of $A$.
\begin{align}
(\lambda^2 - 1)^2 = 0 
&\Rightarrow \lambda^2 = 1\\
&\Rightarrow \lambda = \pm 1
\end{align}
So, $A$ has two distinct eigenvalues: $+1$ and $-1$.

\subsection*{Step 3: Use geometric multiplicity data}

Given:
\begin{align}
\text{GM}(+1) = 3,\quad \text{GM}(-1) = 3.
\end{align}
So total eigenvectors $=6$.  
Since $A$ is $7\times7$, one dimension is missing — corresponding to one generalized eigenvector.

Hence, there must be one Jordan block of size $2$, and five of size $1$.

\subsection*{Step 4: Construct the Jordan form}

Assuming the $2\times2$ block belongs to $\lambda = +1$:
\begin{align}
J =
\myvec{
1 & 1 & 0 & 0 & 0 & 0 & 0\\
0 & 1 & 0 & 0 & 0 & 0 & 0\\
0 & 0 & 1 & 0 & 0 & 0 & 0\\
0 & 0 & 0 & 1 & 0 & 0 & 0\\
0 & 0 & 0 & 0 & -1 & 0 & 0\\
0 & 0 & 0 & 0 & 0 & -1 & 0\\
0 & 0 & 0 & 0 & 0 & 0 & -1
}.
\end{align}

\subsection*{Step 5: Count nonzero entries}

Each $1\times1$ block contributes one nonzero (the eigenvalue).  
The $2\times2$ block contributes three nonzeros (two diagonals and one above it):
\begin{align}
\text{Total nonzeros} = 5\times1 + 3 = 8
\end{align}
\[
\boxed{\text{Answer: 8}}
\]

\section*{Conclusion}

The null spaces $\ker((A-\lambda I)^k)$ reveal how large the Jordan blocks are for each eigenvalue.  
Each difference in their dimensions gives the count of blocks of size $\ge k$.  

Since only one eigenvector is missing from full diagonalizability, exactly one $2\times2$ block exists.  

\[
\boxed{\text{Total nonzero entries in Jordan canonical form of } A = 8.}
\]

\end{document}
